% Editable publishing source. Rebuild without --refresh-tex to preserve edits.
\documentclass[10pt,a4paper]{article}
% Shared editable typography for all five document editions.
\usepackage{fontspec}
\setmainfont{DejaVu Serif}
\setsansfont{DejaVu Sans}
\setmonofont{DejaVu Sans Mono}[Scale=.86]
\usepackage{polyglossia}
\setdefaultlanguage{russian}
\setotherlanguage{english}
\usepackage[a4paper,margin=18mm,top=21mm,bottom=21mm,headheight=15pt]{geometry}
\usepackage{xcolor}
\definecolor{Accent}{HTML}{185A86}
\definecolor{Muted}{HTML}{526573}
\usepackage{graphicx}
\usepackage{longtable,booktabs,array,calc}
\usepackage{fvextra}
\usepackage{pdflscape}
\usepackage{titlesec}
\usepackage{fancyhdr}
\usepackage{enumitem}
\usepackage{etoolbox}
\usepackage{needspace}
\usepackage{tocloft}
\usepackage[unicode,colorlinks=true,linkcolor=Accent,urlcolor=Accent,bookmarksnumbered=false]{hyperref}
\usepackage{bookmark}
\usepackage{microtype}
\AtBeginDocument{\raggedright}
\setcounter{secnumdepth}{0}
\setcounter{tocdepth}{1}
\setlength{\parindent}{0pt}
\setlength{\parskip}{4pt plus 1pt}
\linespread{1.0}
\setlength{\emergencystretch}{3em}
\setlength{\tabcolsep}{4pt}
\renewcommand{\arraystretch}{1.16}
\AtBeginEnvironment{longtable}{\small}
% Measure a compact table or its header and first two data rows. The output
% remains a breakable longtable. Cap the reservation for unusually tall rows.
\newsavebox{\DocumentTableStart}
\newlength{\DocumentTableSpace}
\newcommand{\KeepTableStart}[2]{%
  \setbox\DocumentTableStart=\vbox{\hsize=\linewidth\small #2}%
  \setlength{\DocumentTableSpace}{\dimexpr\ht\DocumentTableStart+\dp\DocumentTableStart
    +\LTpre+\baselineskip+#1\baselineskip*4\relax}%
  \ifdim\DocumentTableSpace>.5\textheight
    \setlength{\DocumentTableSpace}{.5\textheight}%
  \fi
  \Needspace{\DocumentTableSpace}%
}
% Measure short lists once. Small blocks move together; long ones are unboxed
% into the normal page builder and remain breakable, never scaled or clipped.
\newsavebox{\DocumentShortBlock}
\newcommand{\KeepShortBlock}[1]{%
  % A heading's page-space check must not run inside the measuring box: it
  % would inspect the outer page and embed a spurious forced break in the box.
  \par\setbox\DocumentShortBlock=\vbox{\hsize=\linewidth
    \renewcommand{\Needspace}[1]{}#1\par}%
  \ifdim\dimexpr\ht\DocumentShortBlock+\dp\DocumentShortBlock\relax<.32\textheight
    \Needspace{\dimexpr\ht\DocumentShortBlock+\dp\DocumentShortBlock+\medskipamount\relax}%
  \else
    \Needspace{5\baselineskip}%
  \fi
  % Unboxing does not add normal inter-paragraph baseline glue at the leading
  % edge. Give the measured block standard vertical separation from prior text.
  \addvspace{\medskipamount}%
  \unvbox\DocumentShortBlock
}
\preto\section{\Needspace{5\baselineskip}}
\preto\subsection{\Needspace{5\baselineskip}}
\preto\subsubsection{\Needspace{4\baselineskip}}
\clubpenalty=10000
\widowpenalty=10000
\setlist{leftmargin=1.5em,itemsep=2pt,topsep=3pt}
\providecommand{\tightlist}{\setlength{\itemsep}{1pt}\setlength{\parskip}{0pt}}
\titleformat{\section}{\large\sffamily\bfseries\color{Accent}}{}{0pt}{}
\titleformat{\subsection}{\normalsize\sffamily\bfseries\color{Accent}}{}{0pt}{}
\titleformat{\subsubsection}{\normalsize\sffamily\bfseries\color{Accent}}{}{0pt}{}
\titlespacing*{\section}{0pt}{18pt}{8pt}
\titlespacing*{\subsection}{0pt}{13pt}{6pt}
\titlespacing*{\subsubsection}{0pt}{10pt}{5pt}
\pagestyle{fancy}
\fancyhf{}
\fancyfoot[C]{\small\color{Muted}\thepage}
\renewcommand{\headrulewidth}{.35pt}
\renewcommand{\footrulewidth}{0pt}
% A single compact opening page, with profile-controlled TOC depth. Reserve a
% wider page-number column for the multi-page BRD and example editions.
\setlength{\cftbeforesecskip}{5pt}
\setlength{\cftbeforesubsecskip}{2pt}
\renewcommand{\cftsecfont}{\small\sffamily\bfseries}
\renewcommand{\cftsubsecfont}{\small\sffamily}
\renewcommand{\cftsecpagefont}{\small}
\renewcommand{\cftsubsecpagefont}{\small}
\cftsetpnumwidth{2em}
\cftsetrmarg{3em}
\newcommand{\DocumentTitle}[3]{%
  \thispagestyle{plain}
  {\sffamily\bfseries\color{Accent}\fontsize{23}{29}\selectfont #1\par}
  \vspace{3mm}
  {\sffamily\small\color{Muted}#2\par}
  \ifstrempty{#3}{}{\vspace{4mm}#3\par}%
  \begingroup
  \renewcommand{\contentsname}{Содержание}%
  \tableofcontents
  \endgroup
  \vspace{5mm}
  {\color{Accent}\rule{\linewidth}{.8pt}\par}
  \clearpage
}
% Compact standalone profile: keep the title and body on the same opening page.
% It intentionally has no table of contents or forced title-page break.
\newcommand{\CompactDocumentTitle}[3]{%
  \thispagestyle{plain}
  {\sffamily\bfseries\color{Accent}\fontsize{23}{29}\selectfont #1\par}
  \vspace{3mm}
  {\sffamily\small\color{Muted}#2\par}
  \ifstrempty{#3}{}{\vspace{4mm}#3\par}%
  \vspace{5mm}
  {\color{Accent}\rule{\linewidth}{.8pt}\par}
  \vspace{4mm}
}

\hypersetup{pdftitle={Приложение А.1. Требования к входящим артефактам},pdfauthor={},pdfsubject={Рекомендации к BRD}}
\fancyhead[L]{\small SolutionArchitectAI}
\fancyhead[R]{\small Рекомендации к BRD}
\setcounter{tocdepth}{0}
\begin{document}
\CompactDocumentTitle{Приложение А.1. Требования к входящим артефактам}{Рекомендации к BRD}{}
Для обеспечения высокой точности извлечения требований (FR-01), сквозной трассируемости (FR-09) и корректного анализа изменений при обновлении версий (FR-12) входящие документы с бизнес-требованиями должны соответствовать следующим рекомендациям:

\KeepShortBlock{%

\begin{enumerate}
\def\labelenumi{\arabic{enumi}.}
\item
  Метаданные и версионирование: Каждый документ должен содержать уникальный идентификатор, номер редакции (версии) и дату утверждения. Это позволяет системе фиксировать неизменяемую версию источника при привязке архитектурных решений и корректно вычислять разницу (Diff) при загрузке обновлённого файла.
\item
  Сквозная идентификация элементов: Все ключевые сущности внутри документа (бизнес-цели, роли, функциональные и нефункциональные требования, наборы данных, ограничения) должны иметь уникальные текстовые коды (например, FREQ-01, DATA-01, NFR-01). Это исключает неоднозначность при цитировании и упрощает навигацию архитектора от диаграммы к первоисточнику.
\item
  Чёткое разделение границ ответственности (Scope): В документе должны быть явно разграничены функции, реализуемые внутри проектируемой системы, и функции, остающиеся на стороне внешних систем или пользователей. Это предотвращает ошибки топологии графа («чёрные дыры» и «чудеса» по FR-07).
\item
  Структурированное описание данных: Сведения о передаваемой информации должны быть сведены в единый реестр с указанием отправителя, получателя, состава ключевых атрибутов, класса конфиденциальности (например, ПДн, коммерческая тайна) и периодичности передачи.
\item
  Выделение архитектурно значимых NFR и ограничений: Требования к доступности (SLA), пиковым нагрузкам (RPS), окнам пакетной обработки, шифрованию каналов и запретам на использование определённых технологий должны быть описаны в отдельных разделах с привязкой к конкретным функциям или потокам данных.
\item
  Отсутствие скрытых инструкций: Текст документа должен содержать исключительно описание предметной области и требований; любые директивы, имитирующие системные команды для ИИ, игнорируются системой безопасности.
\end{enumerate}

}

Эти рекомендации служат эталоном подготовки входных BRD и не запрещают анализ иных документов MVP. Неоднозначные соответствия требований между редакциями при анализе изменений подтверждает архитектор.

\end{document}
